\section{Discussion and Lessons Learned}
\label{sec:discussion}

\subsection{Treating Writing as Document Engineering}

The central insight from building \software{} is that AI-assisted academic
writing benefits fundamentally from being treated as a document engineering
problem rather than a text generation problem. When the system understands
manuscript structure---chapter boundaries, citation graphs, compilation
dependencies, and file hierarchies---it can provide context-aware assistance
that general-purpose chat interfaces cannot:

\begin{itemize}
  \item \textbf{Automatic context injection:} The system injects relevant
  chapter content and bibliography into AI context without manual copy-paste.
  A single ``review this chapter'' request automatically includes the chapter
  source, the abstract, and related chapter summaries.
  \item \textbf{Verification-driven generation:} Generated citations can be
  verified against scholarly databases in real-time, preventing hallucinated
  references from entering the manuscript.
  \item \textbf{Compilation-aware editing:} Edits can be validated by
  attempting compilation and surfacing errors. This catches missing packages,
  undefined references, and broken cross-references.
\end{itemize}

This engineering mindset also affects how we think about AI quality. A
generated paragraph that reads fluently but cites a non-existent paper or
introduces an undefined \LaTeX{} label is a document engineering failure,
regardless of its textual quality. By centering the system around project
structure rather than text streams, \software{} makes these failures visible
and actionable.

\subsection{The Value of Permission Boundaries}

The three-mode permission model (Chat, Agent, Tools) introduced in
Section~\ref{sec:design} proved to be among the more consequential design
decisions. Based on feedback from five early adopters
(Table~\ref{tab:early-users}), each mode creates a distinct mental model
for the user:

\begin{itemize}
  \item \textbf{Chat mode} became the ``safe space'' for brainstorming,
  discussing reviewer comments, or exploring alternative phrasings. Users
  reported that knowing the AI could not make any changes encouraged more
  exploratory and creative use.
  \item \textbf{Agent mode} became the ``proposal desk.'' Users could see
  exactly what would change before committing. The inline diff visualization
  built trust in the AI's suggestions, even as the scope of edits grew.
  \item \textbf{Tools mode} became the ``workshop'' for operations that span
  multiple steps, such as ``find all passive voice constructions across all
  sections and suggest rewrites.''
\end{itemize}

This separation reduced anxiety about unintended modifications and
encouraged more frequent and more ambitious AI
assistance. Users who were reluctant to use AI in ``auto-edit'' systems
reported comfort with the gated approval workflow.

\subsection{Pipelines as Process Documentation}

An unexpected benefit of the pipeline engine is that executed pipelines serve
as process documentation. Each completed pipeline records:
\begin{itemize}
  \item Which AI stages were run and with what parameters.
  \item What human feedback was provided at each checkpoint.
  \item Whether compilation succeeded or failed and what errors occurred.
  \item The final artifact (PDF) and intermediate outputs.
\end{itemize}
This audit trail is valuable for understanding how a manuscript evolved over
time, for onboarding new collaborators, and for research on writing workflows.
In future versions, this could be formalized as an explicit provenance
subsystem exportable as supplementary material.

\subsection{Local-First as a Practical Advantage}

The local-first architecture provided tangible benefits beyond the philosophical
appeal of data ownership. Manuscripts remain ordinary directories that can be
backed up via rsync, versioned with git, and shared via cloud storage or
USB---all without depending on \software{}. Configuration in \texttt{.env}
files means that different projects can use different LLM providers (OpenAI,
Anthropic, local models via Ollama~\cite{ollama}) without changing application code.

The ability to use external tools (git diff, grep, shell scripts) alongside
\software{} reduced the pressure to implement every conceivable feature within
the application. For example, global search-and-replace across multiple files
is more efficiently performed with \texttt{sed} or VS Code's search feature
than through a custom UI.

\subsection{MCP as an Enabling Standard}

The MCP interoperability layer described in Sections~\ref{sec:design}
and~\ref{sec:implementation} validated the decision to adopt an emerging standard
rather than design a custom API. During testing, Claude Desktop and Cursor discovered
and invoked \software{} tools (paper search, citation verification, \LaTeX{}
compilation) through MCP within minutes of configuration (see the concrete use cases
in Section~\ref{sec:implementation}). In the citation-verification use case, MCP
reduced a 30--45 minute manual cross-referencing task to a 12-second single-command
operation. In the compilation-diagnosis use case, MCP replaced a 3--5 minute
edit--compile--debug cycle with an 8-second tool invocation and structured error
report. These results demonstrate that standards-based interoperability can
significantly reduce integration costs and improve researcher productivity for
repetitive document-engineering tasks. We recommend that new scholarly tools consider
MCP adoption as a lightweight path to AI ecosystem integration.

\subsection{Open Challenges for the Community}
\label{sec:open-challenges}

Three broader challenges emerged from our development experience that merit
attention from the research software community:

\begin{itemize}
  \item \textbf{Evaluation methodology for AI-assisted writing tools.}
  Standard metrics for evaluating AI-assisted writing platforms remain absent.
  Existing work focuses on text quality (perplexity, BLEU, ROUGE), but these
  metrics do not capture workflow integration, user trust, or citation accuracy.
  Developing evaluation frameworks that encompass the full document engineering
  lifecycle---from draft to compilable, citation-verified manuscript---is an
  open problem that the community should address.

  \item \textbf{Trust calibration in AI-augmented workflows.}
  How should users calibrate their trust in AI-generated changes? The
  permission mode system helps, but users may still over-rely on AI
  verification for citations or under-rely on AI suggestions for structure.
  Research on human-AI trust dynamics in document engineering---as distinct
  from general-purpose chat---would inform better interface design.

  \item \textbf{Long-document coherence under AI assistance.}
  For manuscripts exceeding 50 pages (e.g., PhD theses or monographs),
  maintaining global coherence across independently generated or revised
  sections remains challenging, even with project-level context injection.
  This problem is likely to grow as AI-assisted writing scales to book-length
  projects.
\end{itemize}

We discuss \software{}-specific limitations and future development plans in
Section~\ref{sec:limitations}.
